\section*{Conclusion}

Ce premier chapitre a donc permis de montrer la richesse de la pensée politique d'Ambroise malgré l'absence d'un traité exclusivement consacré à une réflexion sur les régimes et le pouvoir. Nous avons donc pu procéder à une analyse des concepts républicains et monarchiques chez Ambroise. Une idée majeure intervient dans l'ensemble de cette théorie : l'importance donnée à la concorde de la société et donc à la relation entre le peuple et ceux ayant le pouvoir, peu importe la forme que prend le gouvernement. S'il se questionne évidemment sur l'exercice du pouvoir pour éviter de tomber dans la coercition et la tyrannie, il valorise tout particulièrement les efforts consentis par les citoyens pour parvenir à une société la plus juste possible. Il cherche à pousser le gouvernement vers le meilleur accomplissement possible. C'est pourquoi la pensée ambrosienne trouve des fondements aussi bien dans l'idéal de la \latin{res publica} que dans un réalisme monarchique. Bien qu'ayant pleinement conscience des risques et des dérives possibles liés à l'exercice du pouvoir par un seul, il reste convaincu que le christianisme possède la capacité d'encadrer et de restreindre cette autorité afin qu'elle atteigne un idéal de moralité chrétienne et qu'elle soit donc compatible avec la concorde de la société.

\bigskip

Et cet idéal politique ne peut donc se réaliser que par le soutien des citoyens allant tous dans la même direction. La théorisation d'un gouvernement romain fondé sur l'idéal chrétien exige selon Ambroise un investissement global des détenteurs du pouvoir et du corps social pour rechercher le respect mutuel et l'accomplissement des devoirs civiques. Cette responsabilisation des citoyens est le cœur de la réflexion politique d'Ambroise : elle vise à garantir l'ordre et la justice par le consensus et l'engagement commun au profit de l'intérêt général. Ainsi, le pouvoir, notamment dans le monde romain, perd toute légitimité à user de la violence et de la coercition pour imposer le respect des lois et conserver le pouvoir. Si une confiance mutuelle s'installe entre le souverain et son peuple, alors les citoyens seront prêts à agir pour la réussite du gouvernement et donc de la société.

\bigskip

Comme nous l'avons vu tout au long du chapitre, sa démarche théorique se fait à travers des textes à destination d'un public large avec le \latin{De Officiis} ou l'\latin{Exameron}. L'évêque cherche ainsi à instaurer une réflexion au sein même de la population, créant une dynamique théorique plus large que de simples traités n'ayant aucun impact concret. Comme nous le verrons dans la suite du mémoire et notamment le chapitre 3, Ambroise agit de façon très directe avec les empereurs, mais sa théorie est vraiment abordée dans ses textes dogmatiques et théologiques.

\bigskip

Nous nous sommes donc ici concentrés sur les fondements théoriques du gouvernement romain idéal selon Ambroise, en questionnant la présence inévitable du modèle impérial, la notion de mérite personnel, les droits des citoyens et l'importance du partage du pouvoir. Mais Ambroise est bien loin de n'être qu'un théologien politique théorique, sa pensée concorde avec une dimension plus pratique où il ne s'interroge pas sur l'exercice du pouvoir et le gouvernement dans sa globalité, mais sur les empereurs qu'il côtoie, leurs agissements et leur lien à la foi chrétienne. La suite de notre analyse de la pensée politique ambrosienne se concentrera donc sur ses écrits s'adressant plus directement aux comportements des empereurs et à la façon de les modeler pour qu'ils continuent de répondre aux exigences d'une monarchie chrétienne.
